% =============================================================================
% MANUEL D'UTILISATION — Template UPC Projet Informatique L3
% Compiler avec : pdflatex > biber > pdflatex > pdflatex
% =============================================================================
\documentclass[12pt, a4paper]{article}
\usepackage{upc}

% --- Informations du projet --------------------------------------------------
\renewcommand{\typeDocument}{Manuel d'utilisation}
\renewcommand{\projetNom}{Nom du projet}
\renewcommand{\projetGroupe}{L3Wx}
\renewcommand{\projetEncadrant}{Prénom Nom}
\renewcommand{\projetAuteurs}{Auteur1, Auteur2, Auteur3, Auteur4}
\renewcommand{\projetVersion}{1.0}
\renewcommand{\projetResume}{%
  Ce manuel d'utilisation décrit l'ensemble des fonctionnalités de l'application
  et guide l'utilisateur dans leur prise en main. Il couvre l'inscription,
  la connexion, l'utilisation des fonctionnalités principales et les actions
  réservées aux administrateurs.
}

\begin{document}

\makecoverpage
\makesommaire

% =============================================================================
\section{Introduction}
% =============================================================================

[Décrire brièvement l'application, son objectif et à qui ce manuel est destiné.]

Ce manuel présente l'ensemble des fonctionnalités de l'application [Nom].
Il s'adresse à deux types d'utilisateurs :
\begin{itemize}
  \item \textbf{Utilisateurs standard} : accèdent aux fonctionnalités
    principales de l'application.
  \item \textbf{Administrateurs} : gèrent le contenu, les utilisateurs et
    la configuration de l'application.
\end{itemize}

% =============================================================================
\section{Inscription \& Connexion}
% =============================================================================

\subsection{Création d'un compte}

[Décrire les étapes pour créer un nouveau compte utilisateur.]

Pour créer un compte :
\begin{enumerate}
  \item Accéder à la page d'accueil de l'application.
  \item Cliquer sur le bouton \textbf{Inscription}.
  \item Renseigner les informations requises :
    \begin{itemize}
      \item Nom d'utilisateur
      \item Adresse e-mail (doit être valide)
      \item Mot de passe (minimum 8 caractères)
      \item Confirmation du mot de passe
    \end{itemize}
  \item Cliquer sur \textbf{Créer mon compte}.
\end{enumerate}

\begin{remarque}
  Une fois le compte créé, vous serez automatiquement connecté et redirigé
  vers la page principale.
\end{remarque}

\subsection{Connexion}

[Décrire les étapes pour se connecter à un compte existant.]

Pour se connecter :
\begin{enumerate}
  \item Accéder à la page de connexion.
  \item Saisir votre adresse e-mail et votre mot de passe.
  \item Cliquer sur \textbf{Se connecter}.
\end{enumerate}

% =============================================================================
\section{Mot de passe oublié}
% =============================================================================

[Décrire la procédure de récupération de mot de passe.]

En cas d'oubli de votre mot de passe :
\begin{enumerate}
  \item Sur la page de connexion, cliquer sur \textbf{Mot de passe oublié ?}
  \item Saisir l'adresse e-mail associée à votre compte.
  \item Cliquer sur \textbf{Envoyer le lien de réinitialisation}.
  \item Consulter votre boîte mail et cliquer sur le lien reçu.
  \item Définir un nouveau mot de passe et confirmer.
\end{enumerate}

\begin{remarque}
  Le lien de réinitialisation est valable 24 heures. Si vous ne recevez pas
  l'e-mail, vérifiez votre dossier de courrier indésirable.
\end{remarque}

% =============================================================================
\section{Fonctionnalités utilisateur}
% =============================================================================

[Décrire chaque fonctionnalité accessible à l'utilisateur standard, illustrée
par des captures d'écran si possible.]

\subsection{Fonctionnalité principale}

[Décrire pas à pas comment utiliser la fonctionnalité principale.]

Pour utiliser cette fonctionnalité :
\begin{enumerate}
  \item Accéder à la section [Nom de la section] depuis le menu principal.
  \item [Étape 2]
  \item [Étape 3]
\end{enumerate}

\begin{remarque}
  [Insérer une capture d'écran avec \texttt{\textbackslash includegraphics}.]
\end{remarque}

\subsection{Fonctionnalité secondaire}

[Décrire la fonctionnalité secondaire.]

\subsection{Profil utilisateur}

[Décrire comment accéder et modifier son profil.]

Pour accéder à votre profil :
\begin{enumerate}
  \item Cliquer sur votre nom d'utilisateur en haut à droite.
  \item Sélectionner \textbf{Mon profil}.
  \item Modifier les informations souhaitées et enregistrer.
\end{enumerate}

\subsection{Déconnexion}

Pour se déconnecter, cliquer sur votre nom en haut à droite puis sur
\textbf{Déconnexion}.

% =============================================================================
\section{Interface administrateur}
% =============================================================================

[Décrire les fonctionnalités réservées aux administrateurs.]

\begin{remarque}[Accès administrateur]
  L'interface d'administration est accessible uniquement aux comptes avec le
  rôle \textbf{Administrateur}. La connexion redirige automatiquement vers
  l'interface d'administration.
\end{remarque}

\subsection{Gestion des utilisateurs}

[Décrire les outils de gestion des utilisateurs.]

L'administrateur peut :
\begin{itemize}
  \item Consulter la liste de tous les utilisateurs inscrits.
  \item Modifier le rôle d'un utilisateur (standard / administrateur).
  \item Désactiver ou supprimer un compte.
\end{itemize}

\subsection{Gestion du contenu}

[Décrire comment l'administrateur gère le contenu de l'application.]

\begin{enumerate}
  \item Accéder à la section \textbf{Gestion du contenu} depuis le menu.
  \item Ajouter un nouvel élément via le bouton \textbf{Ajouter}.
  \item Remplir le formulaire et valider.
\end{enumerate}

\subsection{Création d'un administrateur}

[Expliquer comment créer un nouveau compte administrateur.]

Pour promouvoir un utilisateur au rang d'administrateur :
\begin{itemize}
  \item Depuis la liste des utilisateurs, cliquer sur l'utilisateur ciblé.
  \item Modifier son rôle en \textbf{Administrateur}.
  \item Enregistrer les modifications.
\end{itemize}

% =============================================================================
\section{Conclusion}
% =============================================================================

Ce manuel a présenté l'ensemble des fonctionnalités de l'application [Nom].
Pour toute question ou problème, contacter l'équipe de développement à
l'adresse : \href{mailto:contact@exemple.fr}{contact@exemple.fr}.

% =============================================================================
\section{Références}
% =============================================================================

\printbibliography[heading=none]

\end{document}
